% =============================================================================
% 1.1 Background and Motivation
% =============================================================================

\section{Background and Motivation}
\label{sec:background}

Modern data engineering relies on pipeline orchestration frameworks to automate, schedule, and monitor workflows across unpredictable workloads. A common initial deployment strategy is to host such a framework on a single virtual machine (VM) using a Docker-based run launcher, where each pipeline job runs in its own container on the same host. While simpler than Kubernetes, this architecture carries a fundamental limitation: all containers share the host's finite pool of CPU and memory, and the Linux kernel's OOM killer operates at the host level rather than being scoped to individual containers.

As concurrent job execution increases, resource contention arises. Jobs compete for the same resources, performance degrades non-linearly, and under sufficient load the system fails entirely. Research on shared-memory multiprocessor systems demonstrates that this degradation compounds multiplicatively rather than additively as utilisation approaches capacity limits, producing a collapse effect rather than a smooth performance curve \citep{nanda1991}. \citet{arora2023} further establishes that shared resource contention influences task timing in non-deterministic ways, making worst-case execution behaviour extremely difficult to predict.

Dagster is used in this thesis as a representative modern workflow orchestration system. It provides a Kubernetes Run Launcher that dispatches each pipeline run as an isolated pod on a Kubernetes cluster, where each pod receives defined CPU and memory resources. Failures are theoretically contained within individual pods rather than propagating across the system.

While this architectural migration is well-documented at the tool level, there is limited empirical evaluation of its actual effectiveness under controlled workload conditions. Existing research evaluates Kubernetes performance at the infrastructure level---measuring container overhead, pod scheduling latency, and autoscaler response times \citep{casalicchio2019, nguyen2020, choi2021}. Separately, resource contention theory predicts non-linear collapse in shared execution environments \citep{nanda1991, arora2023}. But no study connects these two bodies of evidence at the application-level workflow execution layer to answer the practical question: \emph{at what workload level does migrating from a shared VM to Kubernetes become worth it?}
